% This file is shared verbatim by the SoftwareX supplement and arXiv appendix.
% Numbering is configured by the including document.

\section{Matched CPU and GPU measurements}
\label{sec:supp-crossover}

Table~\ref{tab:supp-crossover} compares matched CPU and GPU wall times for
identical solver configurations. Ratios are CPU wall time divided by GPU wall
time. CPU and GPU measurements alternated, with each run executed in a separate
process on an Intel Xeon Platinum 8462Y+ and one NVIDIA H100. All measurements
used \texttt{Zipper}, \texttt{SquareSingleNode}, $\beta=3$, and
\texttt{max\_states}=128. Search breadth was fixed to isolate device-dependent
contraction work. The branch-and-bound search runs on the host in both arms. The
36- and 128-spin cases used seven repetitions per device. The 2048-spin cases
used two because of their higher computational cost. The summary uses the
minimum wall time for each arm.

\begin{table}[htbp]
\small
\centering
\caption{CPU/GPU wall-clock ratios for matched configurations. A value below
one indicates shorter CPU wall time. n.m. indicates not measured.}
\label{tab:supp-crossover}
\begin{tabular}{@{}lcccc@{}}
\toprule
Spins and structure & $D{=}8$ & $D{=}16$ & $D{=}32$ & $D{=}64$ \\
\midrule
36, dense     & 0.02 & 0.02 & 0.02 & 0.02 \\
128, dense    & 0.14 & 0.21 & 0.36 & 0.46 \\
128, sparse   & n.m. & n.m. & n.m. & 0.68 \\
2048, dense   & 0.50 & 0.64 & 0.90 & n.m. \\
2048, sparse  & n.m. & 0.78 & \textbf{1.45} & n.m. \\
\bottomrule
\end{tabular}
\end{table}

GPU wall time was shorter only for the 2048-spin sparse instance at $D=32$.
Energies agreed in every measured cell. Raw per-run measurements are stored in
\suppfile{crossover_raw.csv}. The summarized CPU and GPU times and their ratios
are stored in \suppfile{crossover.csv}.
\FloatBarrier

\section{Solution-diversity analysis}
\label{sec:supp-diversity}

The diversity experiments used 2500-spin square-lattice instances without local
fields, so every state $\sigma$ and its global spin flip $-\sigma$ have exactly
the same energy. To avoid counting this symmetry-related pair as two valleys,
distances are reported as
\[
    d_{\mathbb{Z}_2}(\sigma,\tau)
    = \min\!\left(d_H(\sigma,\tau),N-d_H(\sigma,\tau)\right),
\]
where $d_H$ is Hamming distance and $N=2500$. The largest possible quotient
distance is $N/2=1250$. The diversity criterion used in the original article is
$N/4=625$.

Of the 50 instance-temperature runs tested at \texttt{max\_states}=1024, 49
remained within one valley. Across those runs, the largest quotient distance
between any returned state and the best returned state was 399. The sole
exception was instance 1 at $\beta=8$. Its best returned energy was 8.12 above
the lowest energy among the five fixed-order $\beta$ runs for that instance. No
returned state was the global spin flip of the best returned state. We interpret
this result as a degraded contraction rather than evidence of an additional
low-energy valley.

Varying the contraction order broadened the solution set. At $\beta=4$,
selecting one state from each of the eight lattice transformations yielded two
or three distinct valleys on four of five instances. Pairwise quotient distances
reached 1249, compared with the maximum possible value of 1250. Across the same
transformations, the relative energy spread ranged from $2.21\times10^{-4}$ to
$5.74\times10^{-4}$.

The per-run distance quantiles are stored in \suppfile{div3.csv} and
\suppfile{div3_extra.csv}. The per-instance cross-transformation measurements
are in \suppfile{xtransform.csv}. Panel (d) of Figure~1 in the main article
reports medians over ten instances of the 50th, 90th, and 100th within-run
distance percentiles. At $\beta=4$, the square marks the median across five
transformation-sweep instances of the pairwise median distance. Its error bars
extend to the corresponding medians of the pairwise minimum and maximum.

\begin{table}[htbp]
\scriptsize
\centering
\caption{Per-run quotient-distance summaries. Each entry gives
$d_{50}/d_{90}/d_{\max}$ for the 1024 states returned in one fixed-order run.}
\label{tab:supp-diversity-runs}
\begin{tabular}{@{}crrrrr@{}}
\toprule
Instance & $\beta=2$ & $\beta=3$ & $\beta=4$ & $\beta=6$ & $\beta=8$ \\
\midrule
1  & 24/27/31 & 11/15/17 & 7/89/105 & 18/143/145 & 684/702/711 \\
2  & 15/24/31 & 9/143/154 & 20/144/156 & 138/153/164 & 129/152/166 \\
3  & 6/13/29 & 11/19/24 & 10/165/180 & 12/41/66 & 12/40/66 \\
4  & 50/78/94 & 257/293/308 & 8/37/44 & 7/12/52 & 8/13/51 \\
5  & 17/48/63 & 46/60/191 & 42/51/189 & 12/51/190 & 8/45/152 \\
6  & 4/10/13 & 7/13/399 & 10/14/17 & 4/8/11 & 67/71/74 \\
7  & 8/13/19 & 8/12/19 & 32/37/42 & 6/10/13 & 323/329/333 \\
8  & 17/39/45 & 14/27/35 & 28/52/65 & 22/33/40 & 12/19/24 \\
9  & 10/27/32 & 9/24/52 & 15/41/48 & 6/46/51 & 349/352/371 \\
10 & 9/16/26 & 18/28/37 & 16/64/71 & 10/17/23 & 10/17/23 \\
\bottomrule
\end{tabular}
\end{table}

\begin{table}[htbp]
\small
\centering
\caption{Per-instance results across eight lattice transformations at
$\beta=4$. The energy spread is the maximum minus minimum returned energy.
Distances are pairwise quotient distances.}
\label{tab:supp-transform-diversity}
\begin{tabular}{@{}crrrrrrr@{}}
\toprule
Instance & Consensus & $\Delta E$ &
$10^4\Delta E/|E_{\mathrm{best}}|$ &
$d_{\min}$ & $d_{\mathrm{med}}$ & $d_{\max}$ & Valleys \\
\midrule
1 & 1/8 & 0.7815 & 4.02 & 19 & 820 & 1214 & 3 \\
2 & 1/8 & 1.1294 & 5.74 & 21 & 450 & 1249 & 2 \\
3 & 2/8 & 0.4531 & 2.30 & 0  & 971 & 1226 & 2 \\
4 & 1/8 & 0.4309 & 2.21 & 3  & 86  & 686  & 2 \\
5 & 1/8 & 0.4998 & 2.56 & 0  & 212 & 355  & 1 \\
\bottomrule
\end{tabular}
\end{table}
\FloatBarrier

\section{Timing protocol and CUDA memory-pool artifact}
\label{sec:supp-timing}

\subsection{Paired measurements}

The concurrent-sweep and inverse-temperature-ladder comparisons use paired
measurements. The two arms are interleaved within each round. The concurrency
driver performs a full \texttt{GC.gc(true)} before each timed section, followed
by \texttt{CUDA.reclaim()} for device measurements. The ladder driver performs
the same host reclamation before each complete cold or warm arm. Reported
speed-ups are medians of the per-round paired ratios.

The ladder benchmark used the 2048-spin sparse instance at $D=16$ with five
rounds and seeds 2 through 6. Each round ran one cold and one warm ladder on
fresh contractors. Their order alternated across rounds. Rung 1 was cold in
both arms. Rungs 2 and 3 reused a retained boundary MPS only in the warm arm.
Table~\ref{tab:supp-ladder} reports the paired timing ratios. Median discarded
weights for the cold arm were $3.203543\times10^{-3}$,
$1.686507\times10^{-4}$, and $1.856007\times10^{-5}$ over the three rungs. The
warm arm matched the first value and recorded zero for rungs 2 and 3 because a
warm start does not perform a truncating factorization. The returned rung
energies were $-3334.080055$, $-3336.773383$, and $-3336.773383$.

\begin{table}[htbp]
\footnotesize
\centering
\caption{Five-round inverse-temperature-ladder benchmark. Ratios are cold-arm
wall time divided by warm-arm wall time. The interval gives the minimum and
maximum paired ratios. Reduction is $1-1/r$, where $r$ is the median ratio.}
\label{tab:supp-ladder}
\begin{tabular}{@{}cclccc@{}}
\toprule
Rung & $\beta$ & Initialization &
Median ratio & Ratio interval & Reduction \\
\midrule
1 & 1.50 & cold in both arms & 1.008 & [0.916, 1.091] & 0.8\% \\
2 & 2.25 & retained boundary MPS & 1.323 & [1.240, 1.440] & 24.4\% \\
3 & 3.00 & retained boundary MPS & 1.341 & [1.273, 1.730] & 25.4\% \\
\midrule
\multicolumn{3}{l}{Complete ladder} &
1.188 & [1.115, 1.349] & 15.8\% \\
\bottomrule
\end{tabular}
\end{table}

\subsection{Timing contamination checks}

A misleading result from a sequential protocol motivated the paired design. We
retained its summary values, but not the raw development log. A concurrent
warm-up left the CUDA memory pool in a different state when the serial arm was
timed.
For identical work, the serial baseline rose from 14.3~s to 38.8~s, a factor of
approximately 2.7. The contaminated baseline implied a speed-up of
$1.68\times$, whereas the corrected paired protocol gave $0.92\times$. All
reported sweep and ladder ratios use the interleaved protocol and reclamation
steps above.

The crossover experiment uses alternating separate processes. Each process
receives an uncontended device, and the summary uses the minimum of its repeated
wall-time measurements. Timing runs were performed on an otherwise idle
machine.

An exploratory benchmark run coincided with the full package test suite and
returned a different energy at one of 30 benchmark points. We excluded that run
because the machine was not idle. The cause of the difference remains unresolved.
\FloatBarrier

\section{Reproducibility inventory}
\label{sec:supp-reproducibility}

The drivers use the \texttt{SpinGlassPEPS.jl} v2.0.1 API. The launcher runs the
local checkout selected by \texttt{PKG}. To reproduce the article values, check
out the v2.0.1 tag before launching the complete suite from the supplementary
\suppdatadir{} directory with

\begin{verbatim}
PKG=../../SpinGlassPEPS.jl GPU=0 SEED=1 ./run_all.sh
\end{verbatim}

The energy/runtime and fixed-order diversity drivers use
\texttt{SEED + instance}, reset the random stream before each solve, and record
the derived seed in every CSV row. Matching parameter settings use the same
random stream independently of loop order.

The launcher instantiates the package environment, runs the measurement drivers,
writes raw and summarized CSV files, and regenerates the figures.
Table~\ref{tab:supp-inventory} maps the reported results to their drivers and
committed records.

\begin{table}[htbp]
\footnotesize
\centering
\caption{Inventory of supplementary drivers and committed result files.}
\label{tab:supp-inventory}
\begin{tabular}{@{}p{0.24\linewidth}p{0.30\linewidth}p{0.38\linewidth}@{}}
\toprule
Result & Driver or processor & Committed records \\
\midrule
Energy, runtime, and diversity
  & \path{sweep50.jl}, \path{spread.jl}, \path{xtransform.jl}
  & \path{sweep50_cpu.csv}, \path{div3.csv},
    \path{div3_extra.csv}, \path{xtransform.csv} \\
CPU/GPU crossover
  & \path{crossover.jl}, \path{summarize_crossover.py}
  & \path{crossover_raw.csv}, \path{crossover.csv} \\
Concurrent transformation sweep
  & \path{concurrency.jl}, \path{summarize_concurrency.py}
  & \path{concurrency_raw.csv}, \path{concurrency.csv} \\
Allocation and profiling
  & \path{alloc.jl}, \path{profile.jl}
  & \path{alloc_raw.csv}, \path{alloc.csv}, \path{profile.csv} \\
Discarded-weight diagnostics
  & \path{eps_stats.jl}
  & \path{eps_stats.csv} \\
Inverse-temperature ladder
  & \path{ladder.jl}
  & \path{ladder_raw.csv} \\
Figure generation
  & \path{makefigs.py}
  & \path{figures/quality.pdf}, \path{figures/crossover.pdf},
    \path{figures/allocation.pdf} \\
\bottomrule
\end{tabular}
\end{table}

The 36-spin, 128-spin, and 2048-spin instances are distributed in
\path{test/engine/instances/} of the package. The recovered 2500-spin instances
are in \path{benchmark/instances/square_50x50/}. Crossover, concurrency,
post-change allocation, profiling, and ladder measurements used an Intel Xeon
Platinum 8462Y+ system. GPU measurements used one NVIDIA H100.
A multithreaded run reproduced \path{xtransform.csv} byte for byte. Its output
is stored as \path{xtransform_multithread.csv}. The discarded-weight calculation
recorded in \path{eps_stats.csv} ran on the CPU.

\subsection{Discarded-weight record}

Table~\ref{tab:supp-eps} gives the complete diagnostic record used for the
128-spin example in the main article. A truncation is counted as bond-limited
only when the bond cap discards relative weight above $10^{-12}$.

\begin{table}[htbp]
\scriptsize
\centering
\caption{Discarded-weight diagnostics for the 128-spin test instance.}
\label{tab:supp-eps}
\begin{tabular}{@{}crrrrr@{}}
\toprule
$D$ & $\sum_i\varepsilon_i$ & $\max_i\varepsilon_i$ &
Bond-limited & Retained/offered & Energy \\
\midrule
4  & $3.148221\times10^{-4}$ & $2.425283\times10^{-4}$ &
18/18 & 108/307 & $-210.933334$ \\
32 & $5.564114\times10^{-14}$ & $5.511773\times10^{-14}$ &
0/4 & 192/592 & $-210.933334$ \\
\bottomrule
\end{tabular}
\end{table}

\subsection{Allocation provenance}

The pre-change allocation values in Figure~3 of the main article measure
superseded code and cannot be regenerated from v2.0.0. They are retained in
\suppfile{alloc.csv}. Their development code state and protocol are documented
in \suppfile{README.md}. The supplied driver regenerates the final combined
v2.0.0 totals in \suppfile{alloc_raw.csv}.
Table~\ref{tab:supp-allocation} reports the committed isolated-change records in
\suppfile{alloc.csv}. \suppfile{run_all.sh} does not regenerate every
post-change arm in that file.

\begin{table}[htbp]
\small
\centering
\caption{Allocated-memory measurements used in Figure~3 of the main article.}
\label{tab:supp-allocation}
\begin{tabular}{@{}lrrrr@{}}
\toprule
Change & Device & Before [GiB] & After [GiB] & Reduction \\
\midrule
Matrix-backed branched states & GPU & 25.05 & 23.96 & 4.35\% \\
Matrix-backed branched states & CPU & 92.06 & 91.05 & 1.10\% \\
Off-heap temporaries & CPU & 90.9 & 32.0 & 64.80\% \\
\bottomrule
\end{tabular}
\end{table}
\FloatBarrier

\section{Correctness regression coverage}
\label{sec:supp-correctness}

The two result-affecting defects described in the article are covered by
dedicated regression tests. Tests in
\path{test/tensors/contractions_site.jl} and
\path{test/tensors/contractions_virtual.jl} compare
\texttt{corner\_matrix} outputs with independently constructed dense references.
They verify output dimensions and numerical values on CPU and GPU.

The exhaustive-search regression in \path{test/exhaustive/ising.jl} uses a
deterministic ten-spin instance. It compares all 1024 GPU state codes and their
energies with complete CPU enumeration. The GPU launch uses 512 threads per
block, so this test exercises two blocks and crosses the boundary that exposed
the former incomplete-spectrum defect. GPU regression tests run when CUDA is
functional. CPU test groups remain available on systems without a CUDA-capable
device.
